\section{Introduction}

Systolic arrays are a dominant compute substrate in machine-learning (ML) acceleration on field-programmable gate arrays (FPGAs). Their implementations are typically digital signal processing (DSP)-heavy and structurally regular. However, this regularity does not by itself guarantee good timing. Once mapped to a heterogeneous FPGA fabric, a systolic array is strongly affected by long inter-processing-element (PE) connections, so placement has a direct impact on worst negative slack (\wns) and maximum frequency (\fmax).

This setting exposes a clear mismatch between generic FPGA placement and systolic-array placement. General-purpose placers do not explicitly exploit the structural regularity of the array. Our previous work, centered on the \rsad algorithm, addressed this issue by introducing an order-consistency (OC) structural viewpoint for systolic-array placement and showing that this viewpoint can lead to strong half-perimeter wirelength (\hpwl) results~\cite{rsad2023}. The present work builds on that foundation. Its goal is not to replace \rsad as an \hpwl-oriented predecessor, but to extend the same structural line to the timing-driven setting.

The key difficulty is that timing-driven placement is not a single-metric \hpwl problem. On FPGAs, timing is strongly influenced by the longest critical wires, so improving \fmax requires more than reducing aggregate wirelength alone. The method therefore must optimize multiple nets under nonuniform timing-dependent weights while preserving the structural advantages identified by \rsad.

The main contribution of this paper is a formulation result. We first justify the OC condition from the \hpwl viewpoint, and then show that the resulting weighted multi-net objective can be reduced to the optimization of one shortest path on a state DAG. This converts the strip-placement problem into a dynamic-programming problem with exact optimal substructure for fixed weights. It also explains why the \rsad line is scientifically meaningful: once order consistency is justified, the single-column placement problem admits exact optimization rather than only heuristic search.

Based on this formulation, we build a timing-driven placement flow that iteratively adjusts net weights according to timing pressure and repeatedly invokes the DP solver. On the seven benchmarks in our current study, the resulting flow improves average \fmax by 7.4\% over Vivado and by 4.8\% over our previous \rsad-based method. At the same time, it reduces average \hpwl by 24.0\% relative to Vivado, while increasing \hpwl by only 1.9\% relative to \rsad. It also improves average \wns and total negative slack (\tns) over both baselines. These numbers support the central claim of this paper: explicitly controlling the longest wirelength is an effective path toward reducing \wns and improving \fmax for systolic-array DSP placement, while preserving nearly the same \hpwl quality as \rsad.

The contributions of this paper are summarized as follows.
\begin{itemize}[leftmargin=*,topsep=2pt,itemsep=2pt]
\item We further clarify the role of the OC condition introduced by \rsad and show how it turns the single-column DSP placement objective into an ordered form that is compatible with exact shortest-path optimization.
\item We formulate timing-driven DSP placement for systolic arrays on FPGAs through a weighted multi-net objective and show that it can be converted into the shortest-path optimization of a state DAG.
\item We develop a DP solver that is applicable to different net-weight settings, which makes it suitable for timing-driven reweighting rather than only one fixed deterministic placement policy.
\item We build a timing-driven placement flow that improves average \fmax by 7.4\% over Vivado and 4.8\% over the previous \rsad-based method, while reducing average \hpwl by 24.0\% versus Vivado and incurring only a 1.9\% \hpwl increase over \rsad.
\end{itemize}

The rest of the paper is organized as follows. Section~\ref{sec:rw} reviews the most relevant background. Section~\ref{sec:problem} introduces the problem setting and notation. Section~\ref{sec:method} presents the multi-column decomposition, the state-DAG formulation, and the timing-driven flow. Section~\ref{sec:exp} describes the experimental setup and results. Section~\ref{sec:conclusion} concludes the paper.
